\section{Changing one adjacent order}\label{changing-one-adjacent-order}

Let \((B;a,b,c)\) be a rank-three string C-group of type \(\{h,q\}\), with
\(h,q\) even, and suppose the map

\[\rho:B\longrightarrow B_{12},\qquad a\mapsto1, b\mapsto b, c\mapsto c
\tag{7}\]

is a retraction. For \(k\) a positive multiple of \(q\), take the standard
marked epimorphism \(D_{2k}\to B_{12}=D_{2q}\). The compatible product

\[F=\{(g,d)\in B\times D_{2k}:\rho(g)=\overline d\}\tag{8}\]

has markers \((a,1),(b,\widetilde b),(c,\widetilde c)\). The evenness of
\(h\) is forced by (7); it also ensures that the first pair has unchanged
order \(h\).

\textbf{Theorem AL.4 (one-edge extension).} These markers generate \(F\). If \(B\) is finite, its order
is \(|B|k/q\). Its type is \(\{h,k\}\), and it has the intersection
property. The projection onto \(B\) is injective on the first adjacent
subgroup. If \(B\) is solvable, so is \(F\).

\emph{Proof.} Apply UA.SK to the single support $S=\{1,2\}$, using the given deletion and the marked dihedral epimorphism $D_{2k}\to D_{2q}$. Its restrictions to both reflection subgroups are injective, which verifies every proper local subgroup. Its kernel has order $k/q$. The theorem gives generation of the whole compatible product, its order and IP. UA.SK.Local identifies the second adjacent subgroup with $D_{2k}$; the first has trivial projection kernel by (K2), since it does not contain $S$. Thus its order and adjacent product order remain unchanged. Explicitly the other local coordinate of that product has order two, which divides the even integer $h$. The cyclic projection kernel and the solvability of the base prove solvability. \(\square\)

Theorem AL.4 is a specialization of UA.SK; it is included to make clear that
only the indicated deletion is needed here. More simultaneous edges use
the supported coordinate kernels of that theorem.

There is a useful exact presentation variant. Suppose a complete marked
presentation for \(B\) contains \((bc)^q=1\). Replace that relation by
\((bc)^k=1\), retain the remaining relators, and add
\([a,(bc)^q]=1\). Call the resulting presented group \(H\). If every retained
relator holds in (8), then \(H=F\) as a marked group. Here the local dihedral presentation holds in $H$ by the involution relations and $(bc)^k=1$. Its kernel word $(bc)^q$ generates a normal subgroup: $b,c$ invert it and the new relation makes $a$ centralize it. Killing that word gives exactly the original marked presentation of $B$. These are all four conditions of UA.Recognition, so $H=F$ even when the base is infinite. No finite order comparison is needed.
